\chapter{Progettazione del database}\label{chap:Database}
\section{Esportazione Vecchio Db}
La creazione della nuova base di dati è partita dall'esportazione e studio dell'attuale database Access.
L'esportazione è avvenuta per mezzo di \ac{SSMA},
che ha permesso di esportare solo la struttura esistente in SQL Server.
Una volta ricostruita la struttura, essa ha subito un processo di ristrutturazione
che ha portato ad avere un database più leggibile, manuntenibile e performante.
\section{Denominazione tabelle e campi}
Il primo processo di ristrutturazione è stata la creazione delle
 regole di denominazione per le tabelle e le colonne.
Le regole addotatte sono state:
\begin{itemize}
    \item Campi delle tabelle in minuscolo, separati da underscore, ad esempio: \textit{nome\_campo}.
    \item Ad ogni campo viene apportato il prefisso di appartenenza alla tabella, ad esempio: \textit{prefTabella\_nome\_campo}.
    Questa modifica rende ogni campo univoco e facilmente identificabile in caso di join tra tabelle.
    \item Le chiavi primarie di ogni tabella sono degli indici progressivi, rendendo unica ogni riga.
    Il campo segue il pattern \textit{prefTabella\_id}.
\end{itemize}
Nelle tabelle sono stati aggiunti dei campi di audit,
 come la data e l'utente di creazione e di modifica come la data
  e l'utente che ha effettuato l'operazione.
Infine, sempre come campo comune a tutte le tabelle, 
è stato aggiunto il campo boolean \textit{deleted} per la gestione della cancellazione logica.
Questi campi comuni sono stati pensati per essere di supporto nella gestione dei dati, permettendo
la gestione dei backup e delle modifiche senza cancellare mai fisicamente i dati.
\section{Utilizzo di SSIS nel processo di migrazione}
Il processi successivi di modifica e popolamento del database sono stati gestiti tramite \ac{SSIS}.
Questo strumento permette la creazione, anche mediante interfaccia grafica, di pacchetti di query
eseguiti secondo un flusso preimpostato.
Si è quindi creato un primo flusso di query che in ordine:
\begin{enumerate}
    \item Crea se non esiste il database.
    \item Se presenti le tabelle, le elimina.
    \item Crea le tabelle seguendo un ordine consono alle foreignkey.
    \item Popola le tabelle di lookup.
    \item Popola le tabelle con dati del vecchio database.
    \item Popola le tabelle restanti con dati di test.
\end{enumerate}
Questo flusso ha permesso una continua ristrutturazione del database,
 permettendo di testare le modifiche con facilità.
Nelle fasi più avanzate, al flusso prima citato sono state aggiunte le operazioni per la gestione dei backup,
e per la gestione della sincronizzazione dei dati provenienti da altri database.
\section{Ristrutturazione database}
\subsection{Eliminazione ridondanze e normalizzazione}
La fase di studio dell'attuale database ha portato alla luce diverse ridondanze e dati non normalizzati.
All'interno erano presenti più tabella con lo stesso contesto di utilizzo o tabelle non utilizzate.
Si è quindi proceduto all'eliminazione di tali tabelle ed alla normalizzazione delle restanti in seconda forma normale.
Un altra modifica sostanziale è stata l'unificazione di tabelle similari in una unica tabella, discriminate mediante campo tipo proveveniente da un lookup.
Un esempio di questo contesto sono le tabelle \textit{Offerta} e \textit{Ordine}, dove una è la diretta evoluzione della precedente.
Le due tabelle sono state unite sotto il cappello della tabella \textit{Documenti}, discriminata da un campo lookup che ne identifica il tipo.
In queste tabelle in particolare era importante tenere traccia del legame Padre\-Figlio, si è aggiunto quindi un campo \text{doc\_documento_padre} che, mediante
chiave importata, permettessa mediante auto join di ricostruire il legame.
Il metodo così pensato, in ottica di ampliamento, permette la gestione di ulteriori
tipologie di documenti senza cambiare la struttura.
\subsection{Creazione tabelle di lookup}
L'estendibilità del database mediante le dipologie associate alle tabelle passa dalla creazione di tabelle di lookup.
Queste tabelle sono considerate fisse, e non possono essere modificate dall'utente finale,
ma solo da un amministratore tramite query.
Tali tabelle con al di più un campo descrizione, rendono la gestione dei dati più semplice,
potendo filtrare su di essi.
Un vantagio, a livello di software, è il poterle richiamare internamente mediante classi di enumerazione,
 rendendo il codice più leggibile e manutenibile.

\subsection{Aggiunta tabelle non presenti}
\subsubsection{Gestione dei prodotti}
Il databse originale gestiva i prodotti come oggetti atomici, senza la possibilità di gestire le loro componenti.
Si è quindi deciso di creare un auto join n ad n della tabella \textit{Prodotti}.
Tale auto join a livello di logico è stato realizzato con la tabella \textit{Prodotti_composizioni}
che aggiunge, oltre alle chiavi
dei prodotti, anche il quantitativo della riga stessa.
Questo sistema, e mediante appositi flag e query, ha permesso una ricerca approfondita dei ricambi di un prodotto.
Questo sistema ed appositi lookup ha inoltre permesso la costruzione del modulo
 \ac{CPQ} per la configurazione dei prodotti.\ref{chap:Configurazione}.
 %Aggiungere immagine
\subsubsection{Tabelle gestione lingue}
Si è deciso di aggiungere delle apposite tabelle per la gestione delle traduzioni di campi del database.
Le tabelle predisposte contengono al loro interno la chiave della tabella di riferimento,
il campo da tradotto e la chiave importata dalla tabella delle \textit{lingue}.
Quest'ultima tabella non era presente nella soluzione di partenza.
La gestione così pensata permette di estendere le traduzioni a più lingue senza modificare la struttura del database.
Si è anche aggiunta una apposita tabella per gestire le traduzioni nelle reportistiche dei documenti.
Questa tabella mappa le tipologie di reportistica richiedibili dal cliente con le apposite lingue concordate.
Capita infatti di avere ad esempio manuali in lingua inglese, ma etichette sui macchinari con la lingua del luogo.
%Aggiungere immagine
\subsubsection{Creazione tabelle di Back Up}
La scelta delle tabelle soggette a back up è stata a lungo pensata.
Si è scelto di gestire solo un ristretto numero di essenziali tabelle,
 come ad esempio le tabelle dei documenti e dei prodotti, delle righe dei documenti o dei prodotti.
La gestione non è stata capillare per evitare di appesantire il database con tabelle di back up non necessarie.
Una volta decise queste ultime, si è proceduto alla creazione di tabelle di back up,
 con la stessa struttura delle tabelle originali.
Il popolamento avviene mediante dei trigger che scattano ad ogni operazione di Update.
Tramite query e il campo di audit \textit{data_modifica} si è quindi in grado di ricostruire
 lo storico dei dati, e di ripristinare eventuali errori.
 \subsubsection{Tabelle di Utilità}
 Nel database sono state incluse delle tabelle di utilità, che non hanno un legame diretto con le altre tabelle,
 ma saranno utilizzate per la gestione e manutenzione del database.
 Queste tabelle sono:
 \begin{itemize}
    \item \textit{Tabelle}, questa tabella mappa per ogni tabella del database il livello di ruolo necessario per ogni operazione CRUD.
    Tramite apposito look up ad ogni utente è associato un ruolo, corrispondente al livello in ordine decrescente per importanza.
    Si ha quindi 0 per l'Admin e il 6 per il visualizzatore.
    \item \textit{Parametri}, questa tabella contiene solo la versione del database. Questo dato permette
    di verificare se presenti aggiornamenti del database locale prima dell'avvio del sistema.
    \item \textit{Query}, questa tabella raccoglie le query necessarie per aggiornare la struttura del database
    locale con quella in remoto.
 \end{itemize}
\subsection{Importazione dei dati e Stored Procedure}\label{section:importdata}
Come da analisi, è stato richiesto di prelevare e sincronizzare alcuni dati da altri database.
Prima della costruzione delle vere e proprie query di sincronizzazione, si sono
studiate le tabelle di provenienza per ricostruire una ambiente similare facilitando la costruzione delle query.
In particolare, dal database Fusion, si è ricostruita la struttura di asseganzione dinamica delle proprietà ad un prodotto.
Tale struttura, ponendo la cardinalità n ad n tra prodotti e proprietà,
 permette l'assegnazione e la creazione di nuove proprietà senza modificare la struttura del database.
 Alla tabella \text{Proprietà} è associata la tabella \textit{Proprietà\_valori} per la gestione dei valori standard.
 %Immagine struttura Fusion
 Una volta ricostruite le strutture si sono create le seguenti Stored Procedure:
 \begin{itemize}
    \item Da \ac{ERP}, import ed aggiornamento anagrafiche prodotto.
    \item Da \ac{ERP}, import delle distinte prodotto.
    \item Da \ac{PLM}, import ed aggiornamento delle proprietà prodotto.
 \end{itemize}
 Tali stored procedure hanno al loro interno dei controlli che agiscono solo
 quando si verifica una discrepanza.
 Una volta in produzione verranno schedulate giornalmente.